% ============================================================================
% Fichier  : partie4/P04_C12_audit_technique.tex
% Titre    : L'Audit Technique de Sécurité
% Auteur   : MINKA MI NGUIDJOÏ Thierry Emmanuel
% Version  : 1.0 -- novembre 2026
% Corrigés : corriges/Corr_C12.tex
% ============================================================================
% Faits du corpus repris SANS MODIFICATION (ACSI, chapitre 16) : quatre
% centres de LiZbiZ (Douala, Yaoundé, Garoua, Bafoussam) ; sous-domaines
% IN.1 à IN.5 et leurs contrôles clés ; cotes de 2026 par centre (Douala
% 4/1/0, Yaoundé 4/1/0, Garoua 0/4/1, Bafoussam 0/3/2 ; ensemble 8/9/3,
% mu 0,40, U 0,15, H 1,46) ; Garoua et Bafoussam installés par un
% prestataire ; Garoua hébergé par un prestataire local ; loi n° 2010/012,
% art. 69 (cité, texte officiel) et 70 (résumé).
% Sources : CIS Benchmarks (profils de niveau 1 et 2) ; CVSS v4.0 (FIRST,
% 2023) ; catalogue des vulnérabilités exploitées de la CISA (présenté sans
% citation) ; NIST SP 800-53A (méthodes).
% Faits nouveaux (reportés dans la bible) : second passage d'AUDIX en
% octobre 2028, sur exports figés remis le 2 octobre (empreintes SHA-256),
% après le plan B1 à B6 ; cotes 2028 : Douala 5/0/0, Yaoundé 4/1/0, Garoua
% 3/1/1, Bafoussam 2/2/1 ; ensemble 14/4/2 ; Telnet actif sur deux
% commutateurs à Bafoussam ; règle « tout vers tout » à Garoua ; six
% vulnérabilités V1 à V6.
% ============================================================================

\chapter{L'Audit Technique de Sécurité}
\label{chap:audit-technique}
\chaptermark{L'audit technique}

\epigraph{%
    « La politique dit ce qui est permis.\\
    La configuration dit ce qui est possible. »%
}{Principe du cours}

\niveauChapitre{Intermédiaire}{%
    Maîtrise du Chapitre~\ref{chap:smsi}%
}

\begin{boxobjectifs}
\begin{itemize}
    \item Distinguer audit de configuration, scan de vulnérabilités et test
          d'intrusion, et encadrer juridiquement un test technique.
    \item Travailler sur des exports figés et en garantir l'intégrité et
          l'exhaustivité.
    \item Auditer une configuration au regard d'une référence (CIS
          Benchmarks, configuration interne).
    \item Auditer la segmentation par la comparaison des flux autorisés et
          observés.
    \item Auditer les comptes et la journalisation.
    \item Auditer la gestion des vulnérabilités et prioriser leur
          correction.
    \item Comparer deux passages d'audit par centre, sans moyenne
          trompeuse.
\end{itemize}
\end{boxobjectifs}

\bigskip

% ============================================================================
\section*{Avant de commencer : ce qui est écrit est-il configuré~?}
% ============================================================================

Douala, octobre 2028. Deux ans plus tôt, AUDIX CONSEIL avait audité
l'infrastructure des quatre centres de données de LiZbiZ\footnote{LiZbiZ,
AUDIX CONSEIL, PHARMACAM, les dates et les montants sont fictifs~; les
référentiels et les textes cités sont exacts.}~: la politique interdisait
les protocoles d'administration en clair, imposait le changement des mots de
passe par défaut et séparait les postes des serveurs~; deux centres sur
quatre ne l'appliquaient pratiquement pas. Depuis, LiZbiZ a conduit le plan
de traitement de son scénario de rançongiciel (Chapitre~\ref{chap:iso27005}),
et l'audit contractuel de PHARMACAM a encore relevé un filtrage défaillant
entre centres (Chapitre~\ref{chap:audit}).

M.~ETEME demande un second passage. M.~BEYALA accepte~: auditer n'est pas
conseiller, et son cabinet n'a pas conduit le plan. Mais il pose deux
conditions~: travailler sur des exports figés des configurations, et
obtenir une autorisation écrite avant tout test sur les réseaux.

\medskip

Ce chapitre présente l'audit technique de sécurité~: ce qu'il vérifie, ce
qui le distingue d'un test d'intrusion, et comment il produit des preuves
qui tiennent.

\begin{boxalert}
\textbf{Violation typique.} Conclure qu'un centre est « durci » parce que
la politique de durcissement existe et que l'équipe déclare l'avoir
appliquée.

\smallskip
\textbf{Règle~:} seule la configuration réelle, extraite et figée, établit
l'état d'un équipement. L'auditeur la compare à une référence écrite~;
l'écart est le constat.
\end{boxalert}

\separateur

% ============================================================================
\section{Trois Activités Techniques}
\label{sec:c12-activites}
% ============================================================================

\begin{tableaucours}{L{2.6cm}L{3.4cm}L{3.0cm}L{3.2cm}}{%
    \tableauHeader{Activité} &
    \tableauHeader{Question} &
    \tableauHeader{Méthode} &
    \tableauHeader{Limite}
}
Audit de configuration & Les équipements sont-ils configurés selon la
référence~? & Lecture d'exports, comparaison & Ne dit pas ce qu'un
attaquant ferait des écarts \\
Scan de vulnérabilités & Quelles faiblesses connues présentent les
systèmes~? & Outil, de préférence authentifié & Faux positifs~; ne voit
que le connu \\
Test d'intrusion & Un attaquant peut-il atteindre tel objectif~? &
Exploitation contrôlée & Un chemin trouvé ne dit rien des autres \\
\end{tableaucours}
\vspace{-0.8em}
\begin{center}{\small\itshape Tableau~12.1. Trois activités techniques}\end{center}

L'audit technique combine en général les deux premières~; il peut
s'appuyer sur un test d'intrusion, conduit par d'autres (UE MSN-2S3-SEC03),
comme preuve d'efficacité. Ce qui le distingue d'un test d'intrusion, c'est
son objet~: l'audit vérifie l'\emph{ensemble} des critères d'un périmètre et
conclut sur leur conformité~; le test d'intrusion cherche \emph{un} chemin
vers un objectif et conclut sur son existence.

\subsection{Le droit de tester}

\begin{boxjuris}
\textbf{\loicam{n°~2010/012}, article~69}~\cite{loi2010012}~: \textit{« Est
puni d'un emprisonnement de cinq (05) à dix (10) ans et d'une amende de
10.000.000 (dix millions) à 100.000.000 (cent millions) F CFA ou de l'une
de ces deux peines seulement, celui qui accède sans droit, et en violation
des mesures de sécurité, à l'ensemble ou à une partie d'un réseau de
communications électroniques, d'un système d'information ou d'un équipement
terminal, afin d'obtenir des informations ou des données, en relation avec
un système d'information connecté à un autre système d'information. »}

\medskip
\textbf{Article~70} (résumé)~: est puni celui qui provoque, par saturation,
l'attaque d'une ressource de réseau ou d'un système d'information, dans le
but d'empêcher la réalisation des services attendus.

\medskip
\textbf{Portée.} L'autorisation écrite du responsable du système donne à
l'auditeur le droit d'accéder~: sans elle, un scan ou une tentative de
connexion peut tomber sous l'article~69. Les tests de charge et de déni de
service, qui pourraient produire l'effet visé par l'article~70, sont exclus,
sauf autorisation expresse et environnement dédié. Pour le centre de
Garoua, hébergé par un prestataire local, l'accord écrit de l'hébergeur
s'ajoute à celui de LiZbiZ.
\end{boxjuris}

\begin{boxprocedure}
\textbf{La note d'autorisation d'un test technique.} Elle précise~: le
signataire habilité (et l'hébergeur, pour un système hébergé)~; le
périmètre (adresses, systèmes inclus et exclus)~; les méthodes permises et
interdites~; les fenêtres de test~; un contact joignable et une procédure
d'arrêt immédiat~; la confidentialité, la conservation et la destruction
des résultats (articles~14 et~61, Chapitre~\ref{chap:droit}).
\end{boxprocedure}

% ============================================================================
\section{Travailler sur des Exports Figés}
\label{sec:c12-exports}
% ============================================================================

\subsection{Pourquoi figer}

L'auditeur ne travaille pas sur les équipements en production~: il
travaille sur des \termecle{exports figés}\index{export figé}, c'est-à-dire
des copies des configurations, des listes de comptes et des règles de
filtrage, extraites à une date donnée, sous ses yeux ou selon une procédure
écrite, et scellées par une empreinte.

\begin{itemize}
    \item \textbf{Sécurité}~: l'audit ne risque pas de perturber la
          production.
    \item \textbf{Reproductibilité}~: un autre auditeur peut refaire
          l'analyse sur les mêmes fichiers et vérifier les constats.
    \item \textbf{Opposabilité}~: l'empreinte établit que le fichier analysé
          est celui qui a été extrait, à la date dite.
\end{itemize}

\begin{boxprocedure}
\textbf{Le dossier d'exports.} Pour chaque export, le papier de travail
indique~: l'équipement (nom, adresse, rôle, centre)~; la commande ou la
fonction d'export employée~; la date et l'heure~; la personne qui a extrait
et le témoin~; l'empreinte cryptographique (par exemple SHA-256) calculée
au moment de la remise~; le lieu de conservation. L'auditeur recalcule
l'empreinte avant chaque analyse.
\end{boxprocedure}

\subsection{L'exhaustivité}

Un export figé est fiable sur ce qu'il contient~; il ne dit rien de ce qui
manque. L'auditeur rapproche donc la liste des exports de l'inventaire des
équipements. Un équipement présent dans l'inventaire et absent des
exports est un \NEK{}~; un équipement présent dans les exports et absent
de l'inventaire est un constat sur l'inventaire.

\begin{boxscenario}{Opération LIZBIZ. La remise du 2 octobre 2028}
\textbf{Faits.} La DSI remet 61 exports (pare-feu, commutateurs, serveurs,
annuaire), avec leurs empreintes SHA-256. L'inventaire compte 64
équipements. Les trois manquants sont les serveurs de Garoua gérés par
l'hébergeur, qui n'a pas encore fourni ses exports, et un commutateur de
Bafoussam «~en panne~».

\medskip
\textbf{Analyse} (confiance élevée). Les constats sur 61 équipements
peuvent être établis~; pour les trois autres, l'audit cotera \NEK{} et
le dira. Le commutateur «~en panne~» mérite une vérification sur place~: un
équipement que l'on ne peut pas exporter mais qui reste branché est un
équipement que personne ne surveille.
\end{boxscenario}

% ============================================================================
\section{Auditer la Configuration}
\label{sec:c12-configuration}
% ============================================================================

\subsection{Une référence écrite}

Auditer une configuration, c'est la comparer à une référence. Deux sources
sont courantes~:
\begin{itemize}
    \item les \textbf{CIS Benchmarks}, guides de configuration élaborés par
          consensus pour des centaines de produits (systèmes
          d'exploitation, équipements réseau, bases de données), avec deux
          profils~: le niveau~1, applicable presque partout, et le
          niveau~2, plus strict~\cite{cis2026benchmarks}~;
    \item la \textbf{configuration de référence interne}, qui adapte ces
          guides au contexte et que la politique de l'organisation rend
          obligatoire.
\end{itemize}
La seconde est le critère opposable~; la première sert à la construire, et
à juger si elle est suffisante.

\subsection{Lire un export}

\begin{tableaucours}{L{4.0cm}L{4.0cm}L{2.2cm}c}{%
    \tableauHeader{Ligne de l'export (simplifiée)} &
    \tableauHeader{Référence} &
    \tableauHeader{Sous-domaine} &
    \tableauHeader{Cote}
}
\texttt{line vty 0 4 / transport input telnet ssh} & Administration en SSH
seulement & IN.3 & \NON{} \\
\texttt{snmp-server community public RO} & Aucune communauté par défaut &
IN.2 & \NON{} \\
\texttt{username admin secret 9 \$9\$\ldots} & Compte nominatif, mot de
passe haché & IN.2 & \OUI{} \\
\texttt{logging host 10.20.0.15} & Journaux vers le collecteur central &
IN.5 & \OUI{} \\
\texttt{ntp server 10.20.0.1} & Horloge synchronisée & IN.5 & \OUI{} \\
\end{tableaucours}
\vspace{-0.8em}
\begin{center}{\small\itshape Tableau~12.2. Extrait de l'export d'un
commutateur de Bafoussam}\end{center}

Deux lignes suffisent à établir deux non-conformités~: Telnet reste
accepté pour l'administration, et une communauté SNMP par défaut permet de
lire la configuration à quiconque atteint l'équipement. La synchronisation
de l'horloge, souvent négligée, est une condition de la journalisation~: des
journaux mal datés ne permettent pas de reconstituer un incident.

% ============================================================================
\section{Auditer la Segmentation}
\label{sec:c12-segmentation}
% ============================================================================

La segmentation se vérifie en comparant deux matrices~: celle des flux que
la politique autorise entre zones, et celle des flux que les règles des
pare-feu laissent réellement passer. L'écart est le constat.

\begin{tableaucours}{L{3.0cm}cccc}{%
    \tableauHeader{De \textbackslash{} vers} &
    \tableauHeader{DMZ} & \tableauHeader{Applicatif} &
    \tableauHeader{Utilisateurs} & \tableauHeader{Autres centres}
}
Internet & 443 & -- & -- & -- \\
DMZ & -- & 8443 & -- & -- \\
Utilisateurs & 443 & 8443 & -- & -- \\
Administration & 22 & 22 & 22 & 22 \\
Applicatif & -- & -- & -- & 1433 (réplication) \\
\textbf{Maintenance éditeur} & -- & \textbf{toutes (observé)} & -- &
\textbf{toutes (observé)} \\
\end{tableaucours}
\vspace{-0.8em}
\begin{center}{\small\itshape Tableau~12.3. Matrice de flux du centre de
Garoua~: ports autorisés, et en gras les flux observés non autorisés}\end{center}

\begin{boxscenario}{Opération LIZBIZ. Une règle de trop}
\textbf{Constatation.} L'export du pare-feu de Garoua contient une règle
qui autorise tout trafic depuis l'adresse de maintenance de l'éditeur vers
toutes les zones du centre et vers les autres centres. Elle a été créée en
2027, «~temporairement~», lors d'une mise à jour de l'application.

\medskip
\textbf{Analyse} (confiance élevée). La règle annule la segmentation pour
l'accès qui sert de point d'entrée au scénario SO1 (Chapitre~\ref{chap:ebios})~:
un attaquant qui compromettrait l'accès de maintenance atteindrait d'emblée
les quatre centres. C'est la cause du \NON{} relevé par PHARMACAM sur le
filtrage entre centres. Une règle «~temporaire~» sans date d'expiration est
une règle permanente.
\end{boxscenario}

% ============================================================================
\section{Auditer les Comptes et la Journalisation}
\label{sec:c12-comptes}
% ============================================================================

\subsection{Les comptes}

L'audit des comptes rapproche trois listes~: les comptes présents dans les
exports (annuaire, équipements, bases de données), les personnes et les
services autorisés, et les dernières connexions. Il recherche quatre
catégories~:
\begin{itemize}
    \item les \textbf{comptes par défaut} encore actifs (IN.2)~;
    \item les \textbf{comptes orphelins}, sans titulaire (départs)~;
    \item les \textbf{comptes dormants}, sans connexion depuis longtemps~;
    \item les \textbf{comptes partagés} ou génériques, qui empêchent
          d'imputer une action.
\end{itemize}
Les comptes de service, utilisés par les applications, méritent une
attention particulière~: ils ont souvent des droits étendus, des mots de
passe qui ne changent jamais, et personne n'en est propriétaire.

\subsection{La journalisation}

Pour la journalisation (IN.5), l'auditeur vérifie quatre choses~: que
chaque équipement envoie ses journaux au collecteur (rapprochement avec
l'inventaire)~; que les horloges sont synchronisées~; que la durée de
conservation respecte la politique et la loi (article~29,
Chapitre~\ref{chap:droit})~; et que quelqu'un lit ou exploite les journaux,
par des alertes ou une revue. Des journaux collectés et jamais lus ne
détectent rien.

% ============================================================================
\section{Auditer la Gestion des Vulnérabilités}
\label{sec:c12-vulnerabilites}
% ============================================================================

\subsection{Le processus avant le scan}

L'auditeur n'audite pas seulement le résultat d'un scan~; il audite le
processus de gestion des vulnérabilités (ISO/IEC~27002, mesure~8.8)~: les
sources de veille, la fréquence des scans, le délai de correction fixé par
la politique selon la gravité, et le suivi des exceptions. Un scan
\emph{authentifié}, qui se connecte aux systèmes avec un compte de lecture,
voit les versions réellement installées~; un scan non authentifié devine à
partir des services exposés, avec plus de faux positifs.

\subsection{Prioriser}

Le système de cotation CVSS, dans sa version~4.0, attribue à chaque
vulnérabilité un score de gravité de 0 à~10~\cite{first2023cvss4}. Le score
seul ne suffit pas à prioriser~: il décrit la vulnérabilité, pas
l'organisation. Trois questions s'ajoutent~: la vulnérabilité est-elle
exploitée activement (l'agence américaine CISA publie un catalogue des
vulnérabilités exploitées connues)~? Le système est-il exposé (Internet,
accès de maintenance)~? Que porte-t-il (valeur métier, Chapitre~\ref{chap:risque})~?

\begin{tableaucours}{cL{3.6cm}ccL{3.4cm}}{%
    \tableauHeader{Réf.} &
    \tableauHeader{Système} &
    \tableauHeader{CVSS} &
    \tableauHeader{Exploitée~?} &
    \tableauHeader{Exposition}
}
V1 & Serveur de maintenance à distance (Garoua) & 8,1 & Oui & Accès de
l'éditeur \\
V2 & Base de données de l'application (Douala) & 9,0 & Non & Interne,
zone applicative \\
V3 & Portail de suivi des transporteurs & 7,5 & Oui & Internet \\
V4 & Imprimantes des agences & 9,8 & Non & Réseau des agences \\
V5 & Serveur de fichiers de Yaoundé & 6,5 & Oui & Interne \\
V6 & Commutateurs de Bafoussam (Telnet) & 7,0 & Non & Réseau
d'administration \\
\end{tableaucours}
\vspace{-0.8em}
\begin{center}{\small\itshape Tableau~12.4. Six vulnérabilités relevées
(scores et indications d'exploitation fictifs)}\end{center}

Trié par score CVSS, le tableau place V4 (des imprimantes) en tête. Trié
par risque, il place V1 et V3 en tête~: exploitées activement, exposées,
et sur le chemin du scénario SO1 pour V1. V4 vient bien plus loin~: une
imprimante d'agence n'est pas une valeur métier, et elle n'est pas
exposée à Internet.

% ============================================================================
\section{Deux Passages, Quatre Centres}
\label{sec:c12-passages}
% ============================================================================

\begin{boxlecture}{L'infrastructure de LiZbiZ, 2026 et 2028}
\begin{tableaucours}{L{2.4cm}ccccccc}{%
    \tableauHeader{Centre} &
    \tableauHeader{IN.1} & \tableauHeader{IN.2} & \tableauHeader{IN.3} &
    \tableauHeader{IN.4} & \tableauHeader{IN.5} &
    \tableauHeader{2028} & \tableauHeader{2026}
}
Douala & \OUI{} & \OUI{} & \OUI{} & \OUI{} & \OUI{} & 5/0/0 & 4/1/0 \\
Yaoundé & \OUI{} & \OUI{} & \OUI{} & \NON{} & \OUI{} & 4/1/0 & 4/1/0 \\
Garoua & \NON{} & \OUI{} & \OUI{} & \NEK{} & \OUI{} & 3/1/1 & 0/4/1 \\
Bafoussam & \NON{} & \OUI{} & \NON{} & \OUI{} & \NEK{} & 2/2/1 & 0/3/2 \\
\textbf{Ensemble} & & & & & & \textbf{14/4/2} & \textbf{8/9/3} \\
\end{tableaucours}

\medskip
\textbf{Constatation.} En 2028~: $\mu = 0{,}70$~; $U = 0{,}10$~;
\[
H = -(0{,}7\log_2 0{,}7 + 0{,}2\log_2 0{,}2 + 0{,}1\log_2 0{,}1)
  \approx 0{,}360 + 0{,}464 + 0{,}332 \approx 1{,}16~\text{bit}.
\]
En 2026~: $\mu = 0{,}40$~; $U = 0{,}15$~; $H \approx 1{,}46$~bit.

\medskip
\textbf{Analyse} (confiance élevée). L'ensemble passe de la non-conformité
claire à la conformité claire, et l'incertitude baisse. Mais la moyenne
cache deux choses~: Yaoundé régresse sur les correctifs (deux correctifs
critiques en retard), et Garoua comme Bafoussam gardent un défaut de
segmentation, le même que celui qui ouvre le scénario SO1. Les \NEK{}
viennent des exports manquants (hébergeur de Garoua, commutateur de
Bafoussam). Le rapport conclut centre par centre, et place la règle de
maintenance de Garoua en tête des recommandations.

\medskip
\hfill\lecture{0{,}70}{0{,}10}{1{,}16}
\end{boxlecture}

\begin{boxobjet}{SO1 · Rançongiciel par l'accès de maintenance (octobre 2028)}
\textbf{Phase « rentrer »~:} second facteur de l'éditeur prouvé (avenant
signé en 2028)~; mais V1, vulnérabilité exploitée, sur le serveur de
maintenance de Garoua.

\textbf{Phase « trouver »~:} la règle «~tout vers tout~» depuis l'adresse de
maintenance annule la segmentation~; Telnet persiste sur deux commutateurs
de Bafoussam.

\textbf{Phase « exploiter »~:} sauvegardes hors ligne restaurées
(Chapitre~\ref{chap:audit}).

\textbf{Cotation~:} vraisemblance V3 maintenue, tant que V1 et la règle de
Garoua ne sont pas corrigées~; gravité ramenée de G4 à G3 par la
récupération prouvée~; criticité $3 \times 3 = 9$, toujours intolérable.

\textbf{Priorité~:} corriger V1 et supprimer la règle de maintenance de
Garoua sous huit jours, puis vérifier par un nouvel export.
\end{boxobjet}

% ============================================================================
\section{Ce que l'Audit Technique demande à chaque Métier}
\label{sec:c12-metiers}
% ============================================================================

\begin{boxmetier}{Auditeur en cybersécurité}
\textbf{Lire un export}
\begin{enumerate}
    \item Vérifier l'empreinte, la date et l'équipement de chaque export.
    \item Rapprocher les exports de l'inventaire, et coter \NEK{} ce qui
          manque.
    \item Comparer chaque configuration à la référence écrite, ligne par
          ligne.
    \item Comparer les flux observés aux flux autorisés.
    \item Prioriser les vulnérabilités par le risque, pas par le seul score.
    \item Conclure par centre, avant toute agrégation.
\end{enumerate}
\end{boxmetier}

\begin{boxmetier}{RSSI}
\textbf{Préparer l'audit technique}

Fournir la matrice des flux autorisés et la configuration de référence~;
faire signer l'autorisation, avec l'accord des hébergeurs~; prévenir les
équipes de surveillance~; produire les exports avec leur empreinte~;
supprimer les règles « temporaires » à leur échéance.
\end{boxmetier}

\begin{boxmetier}{Gestionnaire des risques cyber}
\textbf{Relier les constats aux scénarios}

Chaque constat technique se rattache à une phase d'un scénario. La règle
de Garoua n'est pas un défaut de configuration parmi d'autres~: elle annule
la phase « trouver » de SO1. C'est ce rattachement qui fixe la priorité.
\end{boxmetier}

\separateur

\begin{boxresume}
\textbf{Ce qu'il faut retenir de ce chapitre~:}
\begin{itemize}
    \item \textbf{Audit de configuration, scan, test d'intrusion}~: trois
          activités, trois questions~; l'audit conclut sur un ensemble de
          critères.
    \item \textbf{Le droit de tester} vient d'une autorisation écrite~; sans
          elle, l'article~69 de la loi n°~2010/012 menace l'auditeur.
    \item \textbf{Les exports figés} rendent l'audit sûr, reproductible et
          opposable~; leur exhaustivité se vérifie contre l'inventaire.
    \item \textbf{Une configuration} se compare à une référence écrite~; une
          segmentation, à la matrice des flux autorisés.
    \item \textbf{Les vulnérabilités} se priorisent par l'exploitation,
          l'exposition et la valeur métier, pas par le seul score CVSS.
    \item \textbf{Quatre centres}~: on conclut centre par centre~; la
          moyenne de LiZbiZ progresse, mais le défaut qui ouvre SO1 demeure.
\end{itemize}
\end{boxresume}

\begin{boxinfo}
\textbf{Points de contrôle.} Vous avez compris ce chapitre si vous savez~:
\begin{itemize}
    \item dire ce qui distingue un audit technique d'un test d'intrusion~;
    \item constituer un dossier d'exports figés~;
    \item lire une ligne d'export et la coter au regard d'une référence~;
    \item prioriser des vulnérabilités en justifiant l'ordre.
\end{itemize}
\end{boxinfo}

% ============================================================================
\section*{Exercices}
% ============================================================================

\begin{boxexo}[title={\ding{45}\quad Exercice~12.1. Quelle activité, quelle autorisation~? \niveaubase}]
Pour chacune des demandes suivantes, dites s'il s'agit d'un audit de
configuration, d'un scan de vulnérabilités ou d'un test d'intrusion, et
qui doit signer l'autorisation.
\begin{enumerate}[(a)]
    \item Vérifier que les pare-feu des quatre centres appliquent la
          matrice de flux.
    \item Savoir si un attaquant parti du portail des transporteurs peut
          atteindre la base de facturation.
    \item Recenser les versions vulnérables des serveurs de Garoua,
          hébergés par un prestataire.
    \item Vérifier qu'aucun commutateur n'accepte Telnet.
    \item Mesurer la résistance du portail à un afflux massif de requêtes.
\end{enumerate}
\end{boxexo}

\begin{boxexo}[title={\ding{45}\quad Exercice~12.2. Lire un export \niveaumoyen}]
L'export d'un serveur Linux de Yaoundé contient, entre autres, les lignes
suivantes (simplifiées)~: \texttt{PermitRootLogin yes}~;
\texttt{PasswordAuthentication yes}~; \texttt{Protocol 2}~; un compte
\texttt{backup} sans connexion depuis 400 jours~; aucune ligne d'envoi des
journaux vers le collecteur.
\begin{enumerate}[(a)]
    \item Pour chaque élément, donnez le sous-domaine (IN.1 à IN.5), la
          référence attendue et la cote.
    \item Lequel de ces écarts touche directement la phase « rentrer » du
          scénario SO1~?
    \item La DSI affirme que le serveur a été corrigé la veille. Que faites-vous~?
\end{enumerate}
\end{boxexo}

\begin{boxexo}[title={\ding{45}\quad Exercice~12.3. Comparer les flux \niveaumoyen}]
La politique de la BANQUE SANAGA n'autorise, depuis le réseau des agences
vers le centre de données, que le port 8443 (application de guichet) et le
port 53 (résolution de noms) vers deux serveurs désignés. L'export du
pare-feu montre en outre~: le port 445 vers le serveur de fichiers du
siège~; le port 3389 vers tous les serveurs~; le port 8443 vers
l'application de guichet.
\begin{enumerate}[(a)]
    \item Construisez la matrice et relevez les flux non autorisés.
    \item Pour chacun, décrivez le risque en vous appuyant sur le scénario
          R-AG-01.
    \item Proposez la recommandation et la preuve de sa mise en œuvre.
\end{enumerate}
\end{boxexo}

\begin{boxexo}[title={\ding{45}\quad Exercice~12.4. Prioriser six vulnérabilités \niveauexpert}]
\begin{enumerate}[(a)]
    \item Classez V1 à V6 (Tableau~12.4) par ordre de priorité de
          correction, et justifiez chaque rang par l'exploitation,
          l'exposition et la valeur métier.
    \item La politique de LiZbiZ fixe un délai de correction de 15 jours pour
          un score supérieur à 9, et de 30 jours au-delà de 7. Montrez
          qu'elle conduit à un mauvais ordre, et proposez une règle
          meilleure.
    \item Le second passage cote IN.4 \NON{} à Yaoundé. Rédigez
          l'observation en sept éléments (Chapitre~\ref{chap:conclure}).
    \item Calculez $\mu$, $U$ et $H$ de Garoua et de Bafoussam en 2028, et
          dites ce que la décomposition par centre apporte par rapport à
          l'ensemble.
\end{enumerate}
\end{boxexo}

\bigskip

\begin{boxinfo}
\textbf{Pour aller plus loin.}
\begin{itemize}
    \item Les CIS Benchmarks~\cite{cis2026benchmarks}, pour les références
          de configuration.
    \item La spécification CVSS~4.0~\cite{first2023cvss4}.
    \item Le NIST SP~800-53A~\cite{nist2022sp80053a}, pour les méthodes
          d'évaluation.
\end{itemize}
\end{boxinfo}

\begin{boxinfo}
\textbf{Transition.} Les audits produisent des constats~; la gestion des
risques produit des scénarios~; la conformité produit des obligations. Il
reste à les faire vivre ensemble, au quotidien, dans un dispositif unique.
C'est la GRC en pratique, objet du Chapitre~\ref{chap:grc-pratique}.
\end{boxinfo}
